一根均匀细直梁长度为 \(L>0\)、抗弯刚度为 \(EI>0\)、单位长度质量为 \(\rho A>0\)，两端均为简支（铰支且无端弯矩）。梁受到大小为 \(P\ge0\) 的恒定轴向压力并保持直梁基态；忽略材料阻尼、非线性挠度及剪切变形。其小横向位移 \(y(x,t)\) 满足
\[
\rho A\,y_{tt}+EI\,y_{xxxx}+P\,y_{xx}=0,
\quad y(0,t)=y(L,t)=0,\quad y_{xx}(0,t)=y_{xx}(L,t)=0.
\]
(1) 用变量分离推导全部正交弯曲模态及相应 \(\omega_n^2(P)\)；(2) 推出稳定直梁的严格负载区间与第一 Euler 临界压力；(3) 写出线性振动能量泛函并证明该临界条件与势能二次型失去正定性一致；(4) 对 \(P\uparrow P_{\rm cr}\) 求最低模频率软化的标度，讨论 \(P>P_{\rm cr}\) 时线性解的增长和为何不能由线性模型预测最终屈曲形状。
